% One-dimensional U-shaped curve. Points and graph edges are schematic.
\begin{tikzpicture}[x=1.55cm,y=1.4cm,font=\footnotesize]
  \draw[BookMuted!55,line width=2.5pt]
    (-1,1.6)--(-1,0) arc[start angle=180,end angle=360,radius=1]--(1,1.6);
  \draw[BookTeal,thick]
    (-1,1.4)--(-1,1.05)--(-1,.7)--(-1,.35)--(-1,0)
    --(-.866,-.5)--(-.5,-.866)--(0,-1)
    --(.5,-.866)--(.866,-.5)--(1,0)--(1,.35)--(1,.7)--(1,1.05)--(1,1.4);
  \foreach \x/\y in {-1/1.4,-1/1.05,-1/.7,-1/.35,-1/0,-.866/-.5,
      -.5/-.866,0/-1,.5/-.866,.866/-.5,1/0,1/.35,1/.7,1/1.05,1/1.4}{
    \fill[FigureInput] (\x,\y) circle (1.7pt);
  }
  \draw[FigureLoss,thick,dashed] (-1,1.4)--(1,1.4)
    node[midway,above] {欧氏捷径};
  \node[left] at (-1,1.4) {$\bm{x}_i$};
  \node[right] at (1,1.4) {$\bm{x}_j$};
  \draw[-{Stealth},BookTeal,thick] (-1.3,.6)--(-1.3,.05);
  \draw[-{Stealth},BookTeal,thick] (1.3,.05)--(1.3,.6);
  \node[BookTeal,align=left,right] at (1.55,-.25) {小邻域连边\\沿曲线绕行};
  \node[BookMuted,below] at (0,-1.2) {灰色曲线：连续流形\qquad 实线折线：图路径};
\end{tikzpicture}
